%%%%%%%%%%%%%%%%
%%% lev 4 %%%
%%%%%%%%%%%%%%%%

\Chapter{4}

\Verse{1}
{
	\hP{וַיְדַבֵּר יְהֹוָה אֶל מֹשֶׁה לֵּאמֹר׃}
}
{
	\eP{And YHWH spoke to Moses, saying,}
}
{
	Ok so YHWH speaks to Drew, saying,
}

\Verse{2}
{
	\hP{דַּבֵּר אֶל בְּנֵי יִשְׂרָאֵל לֵאמֹר נֶפֶשׁ כִּי תֶחֱטָא בִשְׁגָגָה מִכֹּל מִצְוֺת יְהֹוָה אֲשֶׁר לֹא תֵעָשֶׂינָה וְעָשָׂה מֵאַחַת מֵהֵנָּה׃}
}
{
	\eP{
		“Speak to the children of Israel, saying: A person who
		sins by mistake—of any of YHWH’s commandments that are not
		to be done—and does any one of them:
	}
}
{
	Go tell the people that if you sin by mistake—doing anything you're
	not supposed to—and you do any one of them:
%	\footnote{
%		A person who sins. At the time that I am writing this, it has been
%		almost two thousand years since the second Jerusalem Temple was
%		destroyed and sacrifices ended in Judaism. The two largest religions,
%		Christianity and Islam, also do not perform sacrifices, and so it has
%		been so long since we have seen or done this practice that all but a few
%		people are unclear about what sacrifice was. Frequently they imagine
%		that it meant that people regularly destroyed much of their livestock
%		for no reason. This is wrong. The purpose of sacrifice was to recognize
%		that an animal’s life was sacred. The order of the book of Leviticus
%		itself may have contributed to the misunderstanding of sacrifice. It
%		deals with the sacrifices that go to the priests first (in Leviticus
%		1–7) and does not get to the “food” sacrifices until Leviticus 17. Most
%		sacrificed animals are eaten by the persons who bring the animals to the
%		altar. There is also a second group of sacrifices, which involve the
%		ritual slaughter of an animal that is not for consumption by the person
%		who brings the sacrifice. Rather, the meat is eaten by the priests or is
%		entirely burnt. Such sacrifices relate to guilt and expiation. The
%		introduction of law means the introduction of violations of the law.
%		Besides addressing matters of punishment (compensation, execution), the
%		system requires some means of dealing with individuals’ feelings of
%		guilt—and with public condemnation. This is achieved through the
%		sacrificial system, including sacrifices for unwitting violations, by
%		individuals or the community (Leviticus 4 and 5), and sacrifices for
%		various other feelings on the part of the offerer (Leviticus 1–3). The
%		system that provides grounds for guilt also provides a mechanism for
%		expiation and forgiveness. Sacrifice is the only mechanism for
%		forgiveness in the book of Leviticus. There is no suggestion in
%		Leviticus that repentance alone can bring forgiveness for violations of
%		the laws, no indication that one can appeal to YHWH’s mercy, His grace,
%		or His kindness for atonement. Indeed, the words “repentance” (ûb),
%		“mercy” (ramîm), “grace” (n), and “kindness” (esed) do not occur in
%		Leviticus. Thus the psychological and spiritual state of the community
%		is linked powerfully to this visible, tangible act. And this
%		psychological and spiritual focus, together with the very physical focus
%		of the consumption of meat, places sacrifice in a critical and pervasive
%		role in the community’s life. And it places the priests who alone can
%		perform sacrifice in a correspondingly critical role. And, no less
%		significantly, it concentrates the community’s attention on the single
%		place where sacrifice can be performed: the Tabernacle, or Tent of
%		Meeting. Footnote 4:2. by mistake. This begins a lengthy section on what
%		to do when people sin by mistake: a priest, the entire community, a
%		leader, or an individual. Humans make mistakes. But even a sin that is
%		committed by mistake requires some act of atonement. People still feel
%		guilty when they do harm, even if they meant no harm, and so this
%		provides a mechanism for purging the guilt and putting the act in the
%		past. Now, in the absence of sacrifice, other means of atonement have
%		risen in importance. Notably, the Day of Atonement (Yom Kippur) has
%		become the most sacred and widely observed holiday; whereas, in the
%		Torah, Passover stands out as the first and foremost of the holidays.
%	}
}

\Verse{3}
{
	\hP{אִם הַכֹּהֵן הַמָּשִׁיחַ יֶחֱטָא לְאַשְׁמַת הָעָם וְהִקְרִיב עַל חַטָּאתוֹ אֲשֶׁר חָטָא פַּר בֶּן בָּקָר תָּמִים לַיהֹוָה לְחַטָּאת׃}
}
{
	\eP{
		“If the anointed priest will sin, causing the people’s
		guilt, then he shall bring forward for his sin that he
		committed: a bull of the cattle, unblemished, for YHWH as
		a sin offering.
	}
}
{
	Case 1: If it's the oil-headed priests doing it, then he
	has to get a bull—a perfect one—for YHWH as a fine.\fC{
		REF: the anointed priest. Hebrew hakkhn hammîa. This is the first occurrence
		of the word mîa, meaning “anointed” and commonly translated elsewhere in
		the Tanak as “Messiah.” In the Torah mîa always refers to the high
		priest and not, as it later came to mean, the king. The relationship
		between priest and king was always a complex one in the Tanak and in the
		postbiblical world as well. The judge Gideon declines to be made
		Israel’s first king, but he acts as a priest, which may well have
		involved greater power, distinction, and income (Judg 8:22–27). David
		makes some of his sons priests (2Sam 8:18). The Hasmoneans (Maccabees)
		were initially high priests but later were simultaneously priests and
		kings.
	}
}

\Verse{4}
{
	\hP{וְהֵבִיא אֶת הַפָּר אֶל פֶּתַח אֹהֶל מוֹעֵד לִפְנֵי יְהֹוָה וְסָמַךְ אֶת יָדוֹ עַל רֹאשׁ הַפָּר וְשָׁחַט אֶת הַפָּר לִפְנֵי יְהֹוָה׃}
}
{
	\eP{
		And he shall bring the bull to the entrance of the Tent of
		Meeting in front of YHWH. And he shall lay his hand on the
		bull’s head and slaughter the bull in front of YHWH.
	}
}
{
	He has to bring the bull to the Meet Deparmtent in front of YHWH.

	Hand on its head.

	Slaughter it in front of YHWH.
}

\Verse{5}
{
	\hP{וְלָקַח הַכֹּהֵן הַמָּשִׁיחַ מִדַּם הַפָּר וְהֵבִיא אֹתוֹ אֶל אֹהֶל מוֹעֵד׃}
}
{
	\eP{
		And the anointed priest shall take some of the bull’s
		blood and bring it to the Tent of Meeting.
	}
}
{
	Then the guy who did the sin by accident has to take
	some of its blood and bring it to the Tent of Meeting,
	which is where we already are, so it's a pretty quick
	trip.
}

\Verse{6}
{
	\hP{וְטָבַל הַכֹּהֵן אֶת אֶצְבָּעוֹ בַּדָּם וְהִזָּה מִן הַדָּם שֶׁבַע פְּעָמִים לִפְנֵי יְהֹוָה אֶת פְּנֵי פָּרֹכֶת הַקֹּדֶשׁ׃}
}
{
	\eP{
		And the priest shall dip his finger in the blood and
		sprinkle some of the blood seven times in front of YHWH
		before the pavilion of the Holy.
	}
}
{
	Then the priest (not clear if this is the one who did the sin
	thing by accident or another lower level functionary) dips
	his finger in the blood and sprinkle it around in front of
	YHWH in front of the holy prket. Not just once. Seven times.
%	\footnote{
%		the pavilion of the Holy. The pavilion (prket) is over the Holy of
%		Holies, the inner part of the Tabernacle, not over the Holy, the outer
%		part. So why is it called “the pavilion of the Holy” here As in Lev
%		16:2, “the Holy” can refer generally to the holiness that is inside the
%		Holiest of Holies.
%	}
}

\Verse{7}
{
	\hP{וְנָתַן הַכֹּהֵן מִן הַדָּם עַל קַרְנוֹת מִזְבַּח קְטֹרֶת הַסַּמִּים לִפְנֵי יְהֹוָה אֲשֶׁר בְּאֹהֶל מוֹעֵד וְאֵת כּל דַּם הַפָּר יִשְׁפֹּךְ אֶל יְסוֹד מִזְבַּח הָעֹלָה אֲשֶׁר פֶּתַח אֹהֶל מוֹעֵד׃}
}
{
	\eP{
		And the priest shall put some of the blood on the horns of
		the altar of the incense of fragrances, in front of YHWH,
		which is in the Tent of Meeting. And he shall spill all of
		the bull’s blood at the base of the altar of burnt
		offering, which is at the entrance of the Tent of Meeting.
	}
}
{
	He also puts some blood on the altar's horns\fA{It has horns.}
	and on the table with the air freshener on it, in front of YHWH,
	which fortunately is also in the Tent of Meeting where we've
	already come to twice so we're extremely there.

	Now he spills all the blood at the base of the altar of burnt
	offering, which fortunately is also in the Tent of Meeting where we've
	already come to twice so we're extremely there.
}

\Verse{8}
{
	\hP{וְאֶת כּל חֵלֶב פַּר הַחַטָּאת יָרִים מִמֶּנּוּ אֶת הַחֵלֶב הַמְכַסֶּה עַל הַקֶּרֶב וְאֵת כּל הַחֵלֶב אֲשֶׁר עַל הַקֶּרֶב׃}
}
{
	\eP{
		And he shall take off all the fat of the bull of the sin
		offering from it: the fat that covers the innards and all
		the fat that is on the innards
	}
}
{
	Now take all the fat off, especially the visceral fat that he
	may have built up from alcoholism it developed in the course
	of trying to get through all this P material. That fat is the
	worst. So get it off all the innards. If you see an innard,
	get it off of there.
}

\Verse{9}
{
	\hP{וְאֵת שְׁתֵּי הַכְּלָיֹת וְאֶת הַחֵלֶב אֲשֶׁר עֲלֵיהֶן אֲשֶׁר עַל הַכְּסָלִים וְאֶת הַיֹּתֶרֶת עַל הַכָּבֵד עַל הַכְּלָיוֹת יְסִירֶנָּה׃}
}
{
	\eP{
		and the two kidneys and the fat that is on them, that is
		on the loins; and the lobe on the liver: he shall take it
		away with the kidneys
	}
}
{
	Also get rid of the kidneys and also the fat that's on them
	which we already got rid of one sentence ago, but there
	might still be some on there so make sure to double check.

	Also get the fat off the loins and the liver.
	
	Take it away with the kidneys (no period yet we're still talking)
}

\Verse{10}
{
	\hP{כַּאֲשֶׁר יוּרַם מִשּׁוֹר זֶבַח הַשְּׁלָמִים וְהִקְטִירָם הַכֹּהֵן עַל מִזְבַּח הָעֹלָה׃}
}
{
	\eP{
		when it will be taken off from the ox of the
		peace-offering sacrifice. And the priest shall burn them
		to smoke on the altar of burnt offering.
	}
}
{
	when you're taking it off the animal that we're currently
	sacrificing for peace.\fB{Wait I thought this was a sin
	offering, did I fall asleep or something?}

	Anyways then the priest burns them to smoke on the altar
	of the burnt offering. Make sure to use the altar of the
	burnt offering for that. Here's a mnemonic so you don't
	forget: burning starts with b, and so does burnt offering,
	so that's how you can remember.
}

\Verse{11}
{
	\hP{וְאֶת עוֹר הַפָּר וְאֶת כּל בְּשָׂרוֹ עַל רֹאשׁוֹ וְעַל כְּרָעָיו וְקִרְבּוֹ וּפִרְשׁוֹ׃}
}
{
	\eP{
		And the bull’s skin and all of its meat with its head and
		with its legs and its innards and its dung:
	}
}
{
	Also take the bull's skin off.
	
	And all its meat with its head.
	
	And with its legs and its insides.
	
	And get the poop out too, you filthy animal.
}

\Verse{12}
{
	\hP{וְהוֹצִיא אֶת כּל הַפָּר אֶל מִחוּץ לַמַּחֲנֶה אֶל מָקוֹם טָהוֹר אֶל שֶׁפֶךְ הַדֶּשֶׁן וְשָׂרַף אֹתוֹ עַל עֵצִים בָּאֵשׁ עַל שֶׁפֶךְ הַדֶּשֶׁן יִשָּׂרֵף׃}
}
{
	\eP{
		and he shall bring all of the bull outside the camp, to a
		pure place, to the place where ash is spilled, and burn it
		on wood in fire. It shall be burned at the place where ash
		is spilled.
	}
}
{
	Now that the bull's in lots of pieces and not just one, it's
	easier to carry, though you might need to make multiple trips.
	
	So now move the entire bull outside camp, to a pure place,
	the place where we dump out the leftover campfire stuff.

	He (or you (or whoever)) has to burn it there.
}

\Verse{13}
{
	\hP{וְאִם כּל עֲדַת יִשְׂרָאֵל יִשְׁגּוּ וְנֶעְלַם דָּבָר מֵעֵינֵי הַקָּהָל וְעָשׂוּ אַחַת מִכּל מִצְוֺת יְהֹוָה אֲשֶׁר לֹא תֵעָשֶׂינָה וְאָשֵׁמוּ׃}
}
{
	\eP{
		“And if all the congregation of Israel will make a
		mistake, and something will be hidden from the community’s
		eyes, and they do one of any of YHWH’s commandments that
		are not to be done, and they are guilty,
	}
}
{
	So that covers case 1.

	Now for case 2: If all the congregation of Israel makes a
	mistake (presumably that means anyone, not everyone,
	(precise language is important in law)) so anyways if
	all the people make a mistake, and if there's a secret
	about something (we don't know what it is, it's secret)
	and if they do one of the commandments that's not to
	be done, and if on top of doing it they're also guilty
	of having done it (precise language is important in law),
}

\Verse{14}
{
	\hP{וְנוֹדְעָה הַחַטָּאת אֲשֶׁר חָטְאוּ עָלֶיהָ וְהִקְרִיבוּ הַקָּהָל פַּר בֶּן בָּקָר לְחַטָּאת וְהֵבִיאוּ אֹתוֹ לִפְנֵי אֹהֶל מוֹעֵד׃}
}
{
	\eP{
		and the sin over which they have sinned will become known,
		then the community shall bring forward a bull of the
		cattle as a sin offering, and they shall bring it in front
		of the Tent of Meeting.
	}
}
{
	and if the sin that they sinned about stops being a secret,
	then the community (wait maybe it does mean everyone),
	has to bring forward a bull, make sure it's a bull of the
	cattle and not a bull of the ducks or some bull\redacted{shit}
	like that, and they (the community) has to bring it in front
	of the Tent of Meeting.
}

\Verse{15}
{
	\hP{וְסָמְכוּ זִקְנֵי הָעֵדָה אֶת יְדֵיהֶם עַל רֹאשׁ הַפָּר לִפְנֵי יְהֹוָה וְשָׁחַט אֶת הַפָּר לִפְנֵי יְהֹוָה׃}
}
{
	\eP{
		And the community’s elders shall lay their hands on the
		bull’s head in front of YHWH, and he shall slaughter the
		bull in front of YHWH.
	}
}
{
	Then all the olds have to put their hands on the bull's head\fB{
	Talmudic Addendum: As regards the hand on the head situation,
	the olds should do this all at once if they can all fit, but
	it's ok to go one at a time if there's too many olds in your
	community. For example, Miami beach. In cases like
	that, the olds should form a single file line in front of the
	Tent of Meeting and each put the hand on the head one at a time.
	This isn't some random edge case, Miami beach had more jews
	than Israel until recently.} in front of YHWH, and then he
	(the priest) kills the bull in front of YHWH.\fB{If anyone
	needs directions to help them get from ``in front of YHWH''
	to ``in front of YHWH,'' just come to the Tent of Meeting
	in front of YHWH and we'll let you know.}
}

\Verse{16}
{
	\hP{וְהֵבִיא הַכֹּהֵן הַמָּשִׁיחַ מִדַּם הַפָּר אֶל אֹהֶל מוֹעֵד׃}
}
{
	\eP{
		And the anointed priest shall bring some of the bull’s
		blood to the Tent of Meeting.
	}
}
{
	Then the guy with the oiled head brings the blood to the
	Tent of Meeting, where we already are.
}

\Verse{17}
{
	\hP{וְטָבַל הַכֹּהֵן אֶצְבָּעוֹ מִן הַדָּם וְהִזָּה שֶׁבַע פְּעָמִים לִפְנֵי יְהֹוָה אֵת פְּנֵי הַפָּרֹכֶת׃}
}
{
	\eP{
		And the priest shall dip his finger from the blood and
		sprinkle seven times in front of YHWH before the pavilion.
	}
}
{
	Seven sprinkles of blood just like last time.
}

\Verse{18}
{
	\hP{וּמִן הַדָּם יִתֵּן עַל קַרְנֹת הַמִּזְבֵּחַ אֲשֶׁר לִפְנֵי יְהֹוָה אֲשֶׁר בְּאֹהֶל מוֹעֵד וְאֵת כּל הַדָּם יִשְׁפֹּךְ אֶל יְסוֹד מִזְבַּח הָעֹלָה אֲשֶׁר פֶּתַח אֹהֶל מוֹעֵד׃}
}
{
	\eP{
		And he shall put some of the blood on the horns of the
		altar that is in front of YHWH, that is in the Tent of
		Meeting, and he shall spill all of the blood at the base
		of the altar of burnt offering, which is at the entrance
		of the Tent of Meeting.
	}
}
{
	Blood on the horns too. If anyone needs directions to the altar,
	it's right in front of YHWH, that is, in the Tent of Meeting.
	You can't miss it.\fB{Don't come inside though or we'll kill you.
	Fire will come out from YHWH and consume you and your wives
	and your sons and your sons wives with you. You know the drill.}
}

\Verse{19}
{
	\hP{וְאֵת כּל חֶלְבּוֹ יָרִים מִמֶּנּוּ וְהִקְטִיר הַמִּזְבֵּחָה׃}
}
{
	\eP{
		And he shall take off all of its fat from it and burn it
		to smoke at the altar.
	}
}
{
	Take all the fat off and burn it up.
}

\Verse{20}
{
	\hP{וְעָשָׂה לַפָּר כַּאֲשֶׁר עָשָׂה לְפַר הַחַטָּאת כֵּן יַעֲשֶׂה לּוֹ וְכִפֶּר עֲלֵהֶם הַכֹּהֵן וְנִסְלַח לָהֶם׃}
}
{
	\eP{
		And he shall do to the bull as he did to the bull of the
		sin offering. So he shall do to it. And the priest shall
		make atonement over them, and it will be forgiven for
		them.
	}
}
{
	Just do the same thing as before.
	
	\aB{This verse has a special significance in the theology
	of the book, since it shows that even P is getting tired
	of writing all this.}

	\aC{We're not gonna write all that, just do the thing and
	you're forgiven.}

	Delegating rest of today's work to the Most High who is in the tent of meeting at YHWH's ark.\fA{Lit. The box of \heb{ח(shh)ם}.}
}

\Verse{21}
{
	\hP{וְהוֹצִיא אֶת הַפָּר אֶל מִחוּץ לַמַּחֲנֶה וְשָׂרַף אֹתוֹ כַּאֲשֶׁר שָׂרַף אֵת הַפָּר הָרִאשׁוֹן חַטַּאת הַקָּהָל הוּא׃}
}
{
	\eP{
		And he shall bring the bull outside the camp and burn it
		as he burned the first bull. It is the community’s sin
		offering.
	}
}
{
	Burn it like the first bull.

	This completes the description of case 2.
}

\Verse{22}
{
	\hP{אֲשֶׁר נָשִׂיא יֶחֱטָא וְעָשָׂה אַחַת מִכּל מִצְוֺת יְהֹוָה אֱלֹהָיו אֲשֶׁר לֹא תֵעָשֶׂינָה בִּשְׁגָגָה וְאָשֵׁם׃}
}
{
	\eP{
		“When a chieftain will sin and do one of any of the
		commandments of YHWH, his God, that are not to be done, by
		mistake, and he is guilty,
	}
}
{
	Ok now for case 3: If a chieftain sins, and if said chieftain
	also does one of the commandments of YHWH,\fB{His God.} the
	ones that aren't to be done, by mistake, and if on top of
	doing the sin he's also guilty of doing it,\fB{We're gonna
	be here forever aren't we.} don't put a period yet folks
	the sentence isn't done,
%	\footnote{
%		chieftain. Hebrew nî’. Elsewhere in the Torah the word refers to the
%		chieftain of a tribe (Num 2:3–29) or of a priestly family (Num 3:24–35)
%		or of a city (Gen 34:2). Here it may allude to a king of Israel, “a
%		chieftain among your people” (as in 1 Kings 11:34; Ezek 34:24), or it
%		may mean any of these leaders.
%	}
}

\Verse{23}
{
	\hP{אוֹ הוֹדַע אֵלָיו חַטָּאתוֹ אֲשֶׁר חָטָא בָּהּ וְהֵבִיא אֶת קרְבָּנוֹ שְׂעִיר עִזִּים זָכָר תָּמִים׃}
}
{
	\eP{
		or his sin by which he has sinned has been made known to
		him, then he shall bring his offering, a goat, male,
		unblemished.
	}
}
{
	or\fB{Or?}%
	\fC{
		The precedence of logical operators is unspecified in the
		text, though according to a prominent midrash, it goes:
		\texttt{!} (NOT),
		\texttt{\&\&} (AND),
		\texttt{||} (OR).
	}%
	\fA{
		Note: while there's no XOR operator in Classical Biblical Hebrew,
		some distributions define it as a macro in \texttt{/include/asm/yhw.h}.
		See \texttt{man 7 operator} for details on where to put parentheses
		before an expression is executed after placing one's hand on the
		expression's head and then executing it at the Tent of Meeting
		in front of YHWH before the priest dips his finger in the blood of
		the expression and sprinkles it around the altar seven times
		at the Tent of Meeting in front of YHWH.
		Ok back to this bulls\eR{ituation and t}h\eR{e proper procedures for
		sacrificing }it.%
	}
	if the sin by which he sinned is something he knows about\fB{Kill me.}
	then he should bring his offering, a goat, no girl goats don't be cheap,
	and it has to be perfect too don't be cheap.

	Make a note of that and let's keep moving.
	
	God \eR{is good and he will} help us \eR{if we follow the proper protocols
	as outlined in this handbook. I am YHWH, who brought you out of the land
	of Egypt, from a house of bondage, and who made you into a nation
	and multiplied your seed like the starts of the sky and brought
	you into the land I'm giving you to dispossess it before the
	residents of that land that I am giving you to possess. I am YHWH.}.
}

\Verse{24}
{
	\hP{וְסָמַךְ יָדוֹ עַל רֹאשׁ הַשָּׂעִיר וְשָׁחַט אֹתוֹ בִּמְקוֹם אֲשֶׁר יִשְׁחַט אֶת הָעֹלָה לִפְנֵי יְהֹוָה חַטָּאת הוּא׃}
}
{
	\eP{
		And he shall lay his hand on the goat’s head and slaughter
		it at the place where he would slaughter a burnt offering
		in front of YHWH. It is a sin offering.
	}
}
{
	And he puts his hand on the goat's head and slaughters it
	at the place where he would slaughter a burnt offering
	in front of YHWH. It is a sin offering.\fB{
		There's a level of boredom above which the brain
		starts having fun as a survival mechanism.
	}
}

\Verse{25}
{
	\hP{וְלָקַח הַכֹּהֵן מִדַּם הַחַטָּאת בְּאֶצְבָּעוֹ וְנָתַן עַל קַרְנֹת מִזְבַּח הָעֹלָה וְאֶת דָּמוֹ יִשְׁפֹּךְ אֶל יְסוֹד מִזְבַּח הָעֹלָה׃}
}
{
	\eP{
		And the priest shall take some of the blood of the sin
		offering with his finger and put it on the horns of the
		altar of burnt offering, and he shall spill its blood at
		the base of the altar of burnt offering.
	}
}
{
		And the priest shall take some of the blood of the sin\fB{Nevermind, not there yet.}
		offering with his finger and put it on the horns of the
		altar of burnt offering, and he shall spill its blood at
		the base of the altar of burnt offering.\fB{This is hell. We're in hell.}
}

\Verse{26}
{
	\hP{וְאֶת כּל חֶלְבּוֹ יַקְטִיר הַמִּזְבֵּחָה כְּחֵלֶב זֶבַח הַשְּׁלָמִים וְכִפֶּר עָלָיו הַכֹּהֵן מֵחַטָּאתוֹ וְנִסְלַח לוֹ׃}
}
{
	\eP{
		And he shall burn all its fat to smoke at the altar like
		the fat of the peace-offering sacrifice. And the priest
		shall make atonement over him from his sin, and it will be
		forgiven for him.
	}
}
{
	Then he burns the fat to smoke like he did earlier back in
	the other cow.

	Then the priest makes atonement all over him and he's forgiven.
}

\Verse{27}
{
	\hP{וְאִם נֶפֶשׁ אַחַת תֶּחֱטָא בִשְׁגָגָה מֵעַם הָאָרֶץ בַּעֲשֹׂתָהּ אַחַת מִמִּצְוֺת יְהֹוָה אֲשֶׁר לֹא תֵעָשֶׂינָה וְאָשֵׁם׃}
}
{
	\eP{
		“And if one person from the people of the land will sin by
		mistake by doing one of YHWH’s commandments that are not
		to be done and is guilty,
	}
}
{
	Ok now for case 4: a person.\fB{God is dead and P\eR{eople} killed him.}

	If any \textsl{people} sin by mistake, by doing one of YHWH's commandments
	that aren't to be done, and if they're also guilty,
}

\Verse{28}
{
	\hP{אוֹ הוֹדַע אֵלָיו חַטָּאתוֹ אֲשֶׁר חָטָא וְהֵבִיא קרְבָּנוֹ שְׂעִירַת עִזִּים תְּמִימָה נְקֵבָה עַל חַטָּאתוֹ אֲשֶׁר חָטָא׃}
}
{
	\eP{
		or his sin that he has committed has been made known to
		him, then he shall bring his offering, a goat,
		unblemished, female, for his sin that he committed.
	}
}
{
	or if the person in question knows that he did the thing
	that he did, then he has to bring his offering, a goat,
	no girls, no sh\eR{e goats or little b}itty goats---
	wait scratch that, \textsl{only} girls this time,
	but sh\eR{e goats, while required, still must be
	unblemished, including but not limited to being
	of proper size and therefore neither little nor b}itty
	goats\eR{, they} are still prohibited, for his sin that he committed.
}

\Verse{29}
{
	\hP{וְסָמַךְ אֶת יָדוֹ עַל רֹאשׁ הַחַטָּאת וְשָׁחַט אֶת הַחַטָּאת בִּמְקוֹם הָעֹלָה׃}
}
{
	\eP{
		And he shall lay his hand on the head of the sin offering
		and slaughter the sin offering at the place of burnt
		offering.
	}
}
{
		And he shall lay his hand on the head of the sin offering
		and slaughter the sin offering at the place of burnt
		offering.
}

\Verse{30}
{
	\hP{וְלָקַח הַכֹּהֵן מִדָּמָהּ בְּאֶצְבָּעוֹ וְנָתַן עַל קַרְנֹת מִזְבַּח הָעֹלָה וְאֶת כּל דָּמָהּ יִשְׁפֹּךְ אֶל יְסוֹד הַמִּזְבֵּחַ׃}
}
{
	\eP{
		And the priest shall take some of its blood with his
		finger and put it on the horns of the altar of burnt
		offering, and he shall spill all its blood at the base of
		the altar.
	}
}
{
	Blood on finger.
	
	Finger (and therefore blood) on horns.
	
	Horns on altar of burnt offering.

	Rest of the blood: spill it at the base of the altar.\fB{Where's this altar again, can someone remind me?}
}

\Verse{31}
{
	\hP{וְאֶת כּל חֶלְבָּהּ יָסִיר כַּאֲשֶׁר הוּסַר חֵלֶב מֵעַל זֶבַח הַשְּׁלָמִים וְהִקְטִיר הַכֹּהֵן הַמִּזְבֵּחָה לְרֵיחַ נִיחֹחַ לַיהֹוָה וְכִפֶּר עָלָיו הַכֹּהֵן וְנִסְלַח לוֹ׃}
}
{
	\eP{
		And he shall take away all of its fat as the fat was taken
		away from on the peace-offering sacrifice, and the priest
		shall burn it to smoke at the altar as a pleasant smell to
		YHWH, and the priest shall make atonement over him, and it
		will be forgiven for him.
	}
}
{
	Take the fat away like we did when we killed one of those
	previous cows for peace.

	And the priest burns it to smoke at the altar as a pleasant
	smell to YHWH.

	Then the priest makes atonement all over him and he's forgiven.
}

\Verse{32}
{
	\hP{וְאִם כֶּבֶשׂ יָבִיא קרְבָּנוֹ לְחַטָּאת נְקֵבָה תְמִימָה יְבִיאֶנָּה׃}
}
{
	\eP{
		“And if he will bring a lamb as his offering for a sin
		offering, he shall bring a female, unblemished.
	}
}
{
	And if he brings a lamb as the offering, make it
	a girl this time, a hot one.
}

\Verse{33}
{
	\hP{וְסָמַךְ אֶת יָדוֹ עַל רֹאשׁ הַחַטָּאת וְשָׁחַט אֹתָהּ לְחַטָּאת בִּמְקוֹם אֲשֶׁר יִשְׁחַט אֶת הָעֹלָה׃}
}
{
	\eP{
		And he shall lay his hand on the head of the sin offering
		and slaughter it as a sin offering at the place where he
		would slaughter a burnt offering.
	}
}
{
	And put your hand on her head and slaughter her as a sin
	offering at the place where you would slaughter a burnt offering.\fB{Where's that again? I still need that reminder.}
}

\Verse{34}
{
	\hP{וְלָקַח הַכֹּהֵן מִדַּם הַחַטָּאת בְּאֶצְבָּעוֹ וְנָתַן עַל קַרְנֹת מִזְבַּח הָעֹלָה וְאֶת כּל דָּמָהּ יִשְׁפֹּךְ אֶל יְסוֹד הַמִּזְבֵּחַ׃}
}
{
	\eP{
		And the priest shall take some of the blood of the sin
		offering with his finger and put it on the horns of the
		altar of burnt offering, and he shall spill all its blood
		at the base of the altar.
	}
}
{
		And the priest shall take some of the blood of the sin
		offering with his finger and put it on the horns of the
		altar of burnt offering, and he shall spill all its blood
		at the base of the altar.
}

\Verse{35}
{
	\hP{וְאֶת כּל חֶלְבָּהּ יָסִיר כַּאֲשֶׁר יוּסַר חֵלֶב הַכֶּשֶׂב מִזֶּבַח הַשְּׁלָמִים וְהִקְטִיר הַכֹּהֵן אֹתָם הַמִּזְבֵּחָה עַל אִשֵּׁי יְהֹוָה וְכִפֶּר עָלָיו הַכֹּהֵן עַל חַטָּאתוֹ אֲשֶׁר חָטָא וְנִסְלַח לוֹ׃}
}
{
	\eP{
		And he shall take away all of its fat as the sheep’s fat
		was taken away from the peace-offering sacrifice, and the
		priest shall burn them to smoke at the altar with YHWH’s
		offerings by fire, and the priest shall make atonement
		over him, over his sin that he committed, and it will be
		forgiven for him.
	}
}
{
	And now, finally, the chapter ends

	\aB{Praise the lord!}

	\hphantom{And now, finally, the chapter ends} on a cliffhanger, saying...

	That the guy takes the fat away like he did for the sheep earlier.

	And the priest burns it all to smoke at the altar with YHWH's
	offerings by fire.

	Then the priest makes atonement all over him.

	And he's forgiven.
}
